Il concetto di computazione non deterministica viene introdotto per modellare la possibilità di effettuare
transizioni multiple (in questo caso il non determinismo modella una sorta di computazione parallela) oppure
la possibiltà di scegliere tra più possibili evoluzioni della computazione.
Per ottenere una versione non deterministica del concetto di automa a stati finiti introdotto nella sezione
precedente, bisogna variare la definizione nel punto in cui si stabilisce che $Q$ è il codominio della tabella
delle transizioni, sostituendo $Q$ con l’insieme delle parti di $Q$, ovvero $2^Q$: quindi dando alla funzione T la
possibilità di fornire non un unico stato ma un insieme di stati; in questo modo avremo un modello di calcolo
simile al precedente con una dinamica non deterministica poiché può seguire più computazioni allo stesso
tempo, quindi definiamo:
\begin{defi}[Automa Finito Non-deterministico]
Siano $A$ e $Q$ insiemi finiti. Un \emph{automa finito non-deterministico} di alfabeto $A$ e insieme degli stati $Q$ è una tripla $(T,q_0,F)$ dove
\[
T : Q \times A \longrightarrow 2^Q,
\]
\[
F \subseteq Q,
\]
e $q_0\in Q$ è uno stato particolare, detto \emph{stato iniziale}.
\end{defi}

Come nel caso deterministico, associamo alla nozione di automa finito non-deterministico quella di algoritmo formale che fornisce una semantica (operazionale) al modello di calcolo, mostrando come la computazione viene effettuata dall'automa. Per prima cosa è necessario stabilire lo spazio delle configurazioni.

\begin{defi}[Configurazione per Automa Finito Non-deterministico]
Una configurazione per un automa finito non-deterministico di alfabeto $A$ e insieme degli stati $Q$ è una tripla $(w,p,S)$, dove
\begin{enumerate}
    \item $w\in A^*$ è la parola da analizzare;
    \item $p\in\mathbb{N}$ è il puntatore alla posizione della parola che contiene il simbolo dell'alfabeto da considerare per effettuare la successiva transizione;
    \item $S\in 2^Q$ è un insieme di stati.
\end{enumerate}
L'insieme delle configurazioni di alfabeto $A$ e insieme degli stati $Q$ per un automa finito sarà denotato con
\[
\operatorname{Conf}^{A,Q}_{ND}=A^*\times\mathbb{N}\times 2^Q.
\]
Inoltre introduciamo la seguente notazione, quando $S\subseteq Q$:
\[
T(S,a):=\bigcup_{q\in S}\{T(q,a)\}
\]
\end{defi}
Grazie a tale notazione possiamo definire l'algoritmo associato ad un automa finito non-deterministico.

\begin{defi}[Algoritmo associato ad un Automa Finito Non-deterministico]
L'algoritmo associato ad un automa finito non-deterministico di alfabeto $A$ e insieme degli stati $Q$ è la tupla $(A,\{0,1\}, Conf_{ND}^{A,Q},\varepsilon,\tau,\delta)$  dove le tre funzioni si definscoo come segue :

\begin{itemize}
    \item la \emph{funzione di ingresso}
    \[
    e:A^*\longrightarrow {Conf}^{A,Q}_{ND},
    \]
    definita da
    \[
    e(w)=(w,1,\{q_0\});
    \]

    \item la \emph{funzione di transizione}
    \[
    \tau:{Conf}^{A,Q}_{ND}\longrightarrow
    {Conf}^{A,Q}_{ND},
    \]
    definita da
    \[
    \tau((w,p,S))
    =
    (w,p+1,T(S,w(p)));
    \]

    \item la \emph{funzione di uscita}
    \[
    \delta:{Conf}^{A,Q}_{ND}\longrightarrow\{0,1\}^*,
    \]
    definita da
    \[
    \delta((w,p,S))
    =
    \max_{q\in S}F'(q),
    \qquad
    \text{se }p=|w|+1.
    \]
\end{itemize}
\end{defi}

Di conseguenza resta definito il concetto di decidibilità di un insieme mediante automi finiti non-deterministici.

\begin{defi}[Insieme decidibile per Automa Finito Non-deterministico]
Un insieme $X\subseteq A^*$ è detto \emph{decidibile per automa finito non-deterministico} (o \emph{NDFA-decidibile}) se esistono un insieme finito $Q$ ed un automa finito non-deterministico di alfabeto $A$ e insieme degli stati $Q$ tali che l'algoritmo associato decide $X$.
\end{defi}
Un primo esempio di automa finito non-deterministico:

\begin{exam}[Un esempio di NDFA]
\label{es:14}
 Consideriamo $A=\{0,1\}$, $Q=\{ q_0,q_1,q_2,q_3\}$ e definiamo il seguente  NDFA
\[
(T,q_0,F)
\]
dove $F:=\{q_3\}$ e 
\[
\begin{matrix}
T: & Q \times \{ 0,1\} & \to & 2^Q
   \\
   &(q_0, 0) & \mapsto & \{q_0\} \\
   &(q_0, 1) & \mapsto & \{q_0, q_1\} 
    \\
    &(q_1, 0) &\mapsto& \{q_2\} 
    \\ &
    (q_1, 1) & \mapsto & \emptyset 
    \\&(q_2, 0) & \mapsto & \emptyset & 
    \\&(q_2, 1) 
    & \mapsto &
    \{q_3\} \\
&(q_3, 0) & \mapsto & \{q_3\} 
\\&(q_3, 1) & \mapsto & \{q_3\}
\end{matrix}
\]
che può essere rappresenato graficamente come
\[
\begin{tikzpicture}[shorten >=1pt, node distance=2.5cm, auto,
    every state/.style={minimum size=1cm}]
  \node[state, initial] (q0) {$q_0$};
  \node[state] (q1) [right of=q0] {$q_1$};
  \node[state] (q2) [right of=q1] {$q_2$};
  \node[state, accepting] (q3) [right of=q2] {$q_3$};

  \path[->]
    (q0) edge [loop below] node {$0,1$} (q0)
    (q0) edge node {$1$} (q1)
    (q1) edge node {$0$} (q2)
    (q2) edge node {$1$} (q3)
    (q3) edge [loop above] node {$0$} (q3)
    (q3) edge [loop below] node {$1$} (q3);
\end{tikzpicture}
\]
Analizziamo ora il comportamento dell'automa su alcune parole:

\begin{itemize}
    \item Consideriamo la parola
    \[
    w=1101.
    \]
    Calcoliamo l'insieme degli stati raggiungibili dopo la lettura di ciascun simbolo:
    \[
    \begin{aligned}
    \widehat{T}(q_0,1)&=\{q_0,q_1\},\\
    \widehat{T}(q_0,11)
        &=T(q_0,1)\cup T(q_1,1)
        =\{q_0,q_1\}\cup\emptyset
        =\{q_0,q_1\},\\
    \widehat{T}(q_0,110)
        &=T(q_0,0)\cup T(q_1,0)
        =\{q_0\}\cup\{q_2\}
        =\{q_0,q_2\},\\
    \widehat{T}(q_0,1101)
        &=T(q_0,1)\cup T(q_2,1)
        =\{q_0,q_1\}\cup\{q_3\}
        =\{q_0,q_1,q_3\}.
    \end{aligned}
    \]
    Poiché l'insieme finale contiene lo stato finale $q_3$, la parola
    \[
    1101
    \]
    è accettata.

    \item Consideriamo ora la parola
    \[
    w=1000.
    \]
    Si ha
    \[
    \begin{aligned}
    \widehat{T}(q_0,1)&=\{q_0,q_1\},\\
    \widehat{T}(q_0,10)
        &=T(q_0,0)\cup T(q_1,0)
        =\{q_0\}\cup\{q_2\}
        =\{q_0,q_2\},\\
    \widehat{T}(q_0,100)
        &=T(q_0,0)\cup T(q_2,0)
        =\{q_0\}\cup\emptyset
        =\{q_0\},\\
    \widehat{T}(q_0,1000)
        &=T(q_0,0)
        =\{q_0\}.
    \end{aligned}
    \]
    L'insieme degli stati raggiungibili non contiene alcuno stato finale, quindi la parola
    \[
    1000
    \]
    è rifiutata.
    \item Consideriamo infine la parola
\[
w=1100101.
\]
Si ha
\[
\begin{aligned}
\widehat{T}(q_0,1)&=\{q_0,q_1\},\\
\widehat{T}(q_0,11)
    &=T(q_0,1)\cup T(q_1,1)
    =\{q_0,q_1\}\cup\emptyset
    =\{q_0,q_1\},\\
\widehat{T}(q_0,110)
    &=T(q_0,0)\cup T(q_1,0)
    =\{q_0\}\cup\{q_2\}
    =\{q_0,q_2\},\\
\widehat{T}(q_0,1100)
    &=T(q_0,0)\cup T(q_2,0)
    =\{q_0\}\cup\emptyset
    =\{q_0\},\\
\widehat{T}(q_0,11001)
    &=T(q_0,1)
    =\{q_0,q_1\},\\
\widehat{T}(q_0,110010)
    &=T(q_0,0)\cup T(q_1,0)
    =\{q_0\}\cup\{q_2\}
    =\{q_0,q_2\},\\
\widehat{T}(q_0,1100101)
    &=T(q_0,1)\cup T(q_2,1)
    =\{q_0,q_1\}\cup\{q_3\}
    =\{q_0,q_1,q_3\}.
\end{aligned}
\]
Poiché l'insieme finale contiene lo stato finale $q_3$, la parola
\[
1100101
\]
è accettata.
\end{itemize}
Si può mostrare infatti che questo automa è quello che decide l'insieme delle parole che contengono come sottostringa $101$.
Consideriamo ora lo stesso automa sostituendo la transizione
\[
T(q_0,1)=\{q_0,q_1\}
\]
con
\[
T(q_0,1)=\{q_1\},
\]
lasciando invariate tutte le altre transizioni.
Mostriamo che questo automa non riconosce più tutte le parole contenenti la sottostringa \(101\).
Infatti, la parola
\[
w=1101
\]
contiene la sottostringa \(101\), ma l'automa si comporta nel modo seguente:
\[
\begin{aligned}
\widehat{T}(q_0,\varepsilon)&=\{q_0\},\\
\widehat{T}(q_0,1)&=\{q_1\},\\
\widehat{T}(q_0,11)&=T(q_1,1)=\emptyset,\\
\widehat{T}(q_0,110)&=\bigcup_{q\in\emptyset}T(q,0)=\emptyset,\\
\widehat{T}(q_0,1101)&=\bigcup_{q\in\emptyset}T(q,1)=\emptyset.
\end{aligned}
\]
Poiché la computazione termina con l'insieme vuoto, la parola \(1101\) viene rifiutata.

L'errore consiste nel fatto che l'automa è costretto a scegliere il \emph{primo} simbolo \(1\) come possibile inizio della sottostringa \(101\). Nell'automa non deterministico originale, invece, il ramo che rimane nello stato \(q_0\) permette di ignorare il primo \(1\) e di scegliere il secondo \(1\) come inizio della sottostringa, accettando così correttamente la parola.
\end{exam}
In vista, dell'obiettivo di mostrare che l'insieme $X_1 \cdot X_2$ ottenuto a partire da due insiemi $X_1$ e $X_2$ DFA-decidibili è lui stesso DFA-decidibile dobbiamo poter costruire un automa che ``concatena'' la computazione di un automa con quella del secondo, a questo fine è comodo utilizzare una ulteriore variante della definizione di NDFA che può effettuare transizioni senza fare avanzare il puntatore alla posizione corrente:

\begin{defi}[Automa Finito non-deterministico con $\epsilon$-mosse]
\label{def:15}
Siano $A$ e $Q$ insiemi finiti, un automa finito non-deterministico con $\epsilon$-mosse \emph{di alfabeto $A$ e insieme degli stati $Q$} è una tripla $(T, q_0, F)$ dove
\[
T : Q \times A \cup \{\epsilon\} \to 2^Q
\]
\[
F \subset Q
\]
e $q_0 \in Q$ è uno stato particolare, detto \emph{stato iniziale}.
\end{defi}



Per la descrizione della computazione dovremo introdurre una variante della funzione di transizione che prenda in considerazione le $\epsilon$-mosse, a questo scopo introduciamo la seguente notazione per $S \in 2^Q$:
\[
T^k(S,a) :=
\begin{cases}
S & \text{se } k = 0 \\
T(T^{k-1}(S,a), a) & k > 0
\end{cases}
\]

\begin{defi}[Algoritmo associato ad un automa finito non-deterministico con $\epsilon$-mosse]
\label{def:17}
L'algoritmo associato ad un automa non-deterministico di alfabeto $A$ e insieme degli stati $Q$ viene descritto definendo le tre funzioni:
\begin{itemize}
\item la \textbf{funzione di ingresso} che associa ad un parola $w \in A^*$, la configurazione
\[
\varepsilon(w) = \left(w, 1, \bigcup_{k \geq 0} T^k(\{q_0\}, \epsilon)\right),
\]

\item la \textbf{funzione di transizione} $\tau : \mathrm{Conf}_{ND}^{A,Q} \to \mathrm{Conf}_{ND}^{A,Q}$ è data da
\[
\tau((w,p,S)) = \left(w, p+1, \bigcup_{k \geq 0} T^k(T(S, w(p)), \epsilon)\right),
\]

\item la \textbf{funzione di uscita} $\delta : \mathrm{Conf}_{ND}^{A,Q} \to \{0,1\}$, data da
\[
\delta((w,p,S)) =
\begin{cases}
0 & \text{se } p > |w| \text{ e } F \cap S = \emptyset \\
1 & \text{se } p > |w| \text{ e } F \cap S \neq \emptyset
\end{cases}
\]
\end{itemize}
\end{defi}


\begin{exam}[Un esempio di $\epsilon$-NDFA]
\label{es:16}

Consideriamo l'alfabeto
\[
A=\{0,1\},
\]
l'insieme degli stati
\[
Q=\{q_0,q_1,q_2,q_3\},
\]
e definiamo il seguente $\epsilon$-NDFA
\[
(T,q_0,F),
\]
dove
\[
F=\{q_3\}
\]
e
\[
T:Q\times(A\cup\{\epsilon\})\to 2^Q
\]
è definita da
\[
\begin{matrix}
T(q_0,0)&=&\{q_0\},\\
T(q_0,1)&=&\{q_0,q_1\},\\
T(q_1,0)&=&\{q_2\},\\
T(q_1,1)&=&\emptyset,\\
T(q_1,\epsilon)&=&\{q_2\},\\
T(q_2,0)&=&\emptyset,\\
T(q_2,1)&=&\{q_3\},\\
T(q_3,0)&=&\{q_3\},\\
T(q_3,1)&=&\{q_3\}.
\end{matrix}
\]
L'automa può essere rappresentato graficamente come
\[
\begin{tikzpicture}[shorten >=1pt, node distance=2.5cm, auto,
    every state/.style={minimum size=1cm}]
  \node[state, initial] (q0) {$q_0$};
  \node[state] (q1) [right of=q0] {$q_1$};
  \node[state] (q2) [right of=q1] {$q_2$};
  \node[state, accepting] (q3) [right of=q2] {$q_3$};

  \path[->]
    (q0) edge [loop below] node {$0,1$} (q0)
    (q0) edge node {$1$} (q1)
    (q1) edge node {$0,\epsilon$} (q2)
    (q2) edge node {$1$} (q3)
    (q3) edge [loop above] node {$0$} (q3)
    (q3) edge [loop below] node {$1$} (q3);
\end{tikzpicture}
\]
Per descrivere la computazione consideriamo la funzione di transizione estesa
\[
\widehat{T}:2^Q\times A^*\to 2^Q
\]
definita considerando le $\epsilon$-mosse come transizioni gratuite. In
particolare, per ogni insieme di stati $S\subseteq Q$ poniamo
\[
T^0(S,\epsilon)=S,
\]
e
\[
T^k(S,\epsilon)=T(T^{k-1}(S,\epsilon),\epsilon)
\]
per $k>0$.
La computazione dell'automa consiste quindi nel considerare tutti gli stati
ottenibili dopo un numero finito di $\epsilon$-mosse.

\begin{itemize}

\item Consideriamo la parola
\[
w=101.
\]

Partendo dallo stato iniziale otteniamo
\[
\widehat{T}(q_0,\epsilon)=\{q_0\}.
\]

Dopo aver letto il primo simbolo $1$:
\[
T(\{q_0\},1)=\{q_0,q_1\}.
\]

Applicando successivamente le $\epsilon$-mosse:
\[
T^0(\{q_0,q_1\},\epsilon)=\{q_0,q_1\},
\]
\[
T^1(\{q_0,q_1\},\epsilon)
=
T(\{q_0,q_1\},\epsilon)
=
\{q_2\}.
\]

Quindi dopo il primo simbolo:
\[
\widehat{T}(q_0,1)
=
\{q_0,q_1,q_2\}.
\]

Ora leggiamo il simbolo $0$:
\[
T(\{q_0,q_1,q_2\},0)
=
\{q_0\}\cup\{q_2\}\cup\emptyset
=
\{q_0,q_2\}.
\]

Non sono possibili ulteriori $\epsilon$-mosse, quindi:
\[
\widehat{T}(q_0,10)=\{q_0,q_2\}.
\]

Infine leggiamo $1$:
\[
T(\{q_0,q_2\},1)
=
\{q_0,q_1\}\cup\{q_3\}
=
\{q_0,q_1,q_3\}.
\]

Poiché
\[
q_3\in\widehat{T}(q_0,101),
\]
la parola $101$ è accettata.

\item Consideriamo la parola
\[
w=10.
\]

Dopo il simbolo $1$ abbiamo già visto che
\[
\widehat{T}(q_0,1)=\{q_0,q_1,q_2\}.
\]

Leggendo $0$ otteniamo:
\[
T(\{q_0,q_1,q_2\},0)
=
\{q_0\}\cup\{q_2\}\cup\emptyset
=
\{q_0,q_2\}.
\]

Non essendo presente lo stato finale:
\[
q_3\notin\widehat{T}(q_0,10),
\]
la parola $10$ è rifiutata.

\end{itemize}

\end{exam}

Con la nozione appena introdotta è quindi possibile dimostrare in modo diretto che dati due insiemi NDFA-decidibili, l'insieme ottenuto come concatenazione è pure NDFA-decidibile, ovvero la Proposizione~\ref{prop:9} nel caso non-deterministico, e infatti
\begin{prop}
    Siano $X_1$ e $X_2$ decidibili per automa finito non-deterministico allora anche $X_1\cdot X_2$ è
decidibile per automa finito non-deterministico.
\begin{proof}
Siano i due automi che decidono $X_1 \subset A^*$ e $X_2 \subset A^*$ rispettivamente dati su insiemi degli stati disgiunti $Q_1$ e $Q_2$, ovvero con $Q_1 \cap Q_2 = \emptyset$: $\mathcal{A}_1 = (T', q_0', F')$ e $\mathcal{A}_2 = (T'', q_0'', F'')$.

Costruiamo l'automa che decide $X_1 \cdot X_2$ aggiungendo alle transizioni di $\mathcal{A}_1$ e $\mathcal{A}_2$ quelle che vanno dagli stati accettanti di $\mathcal{A}_1$ allo stato iniziale $q_0''$ del secondo automa, ovvero definiamo l'automa $\mathcal{A}$ con insieme degli stati $Q := Q_1 \cup Q_2$ dove lo stato iniziale è $q_0 := q_0' \in Q_1$ e gli stati finali accettanti sono definiti come quelli di $\mathcal{A}_2$, dunque:
\[
\mathcal{A} := (T, q_0, F)
\]
dove
\[
T(q,a) :=
\begin{cases}
T'(q,a) & \text{se } q \in F' \text{ e } a \neq \epsilon \text{ oppure } q \notin F', \\
T'(q,a) \cup \{q_0''\} & q \in F' \text{ e } a = \epsilon, \\
T''(q,a) & q \in Q_2,
\end{cases}
\]
e l'insieme degli stati finali accettanti è dato da $F := F''$.
\end{proof} 
\end{prop}
\begin{prop}
\label{prop:19}
Sia dato un insieme $X \subseteq A^*$ per un certo alfabeto finito $A$.
Allora
\[
X \text{ è DFA-decidibile}
\iff
X \text{ è NDFA-decidibile}
\iff
X \text{ è }\epsilon\text{-NDFA-decidibile}.
\]

\begin{proof}
Procediamo dimostrando le tre implicazioni.

\begin{itemize}

\item[$(\mathrm{DFA}\Rightarrow\mathrm{NDFA})$]
Sia $(T,q_0,F)$ un DFA. Basta considerare la funzione di transizione come funzione a valori nell'insieme delle parti:
\[
T' : Q\times A\longrightarrow 2^Q,
\]
definita da
\[
T'(q,a):=\{T(q,a)\}.
\]
L'automa così ottenuto è un NDFA che produce esattamente la stessa computazione.

\item[$(\mathrm{NDFA}\Rightarrow\epsilon\text{-}\mathrm{NDFA})$]
Sia $(T,q_0,F)$ un NDFA. Definiamo
\[
T':Q\times(A\cup\{\epsilon\})\longrightarrow 2^Q
\]
ponendo
\[
T'(q,a):=T(q,a),
\qquad
a\in A,
\]
e
\[
T'(q,\epsilon):=\varnothing.
\]
In questo modo l'automa non utilizza mai $\epsilon$-mosse e coincide con quello di partenza.

\item[$(\epsilon\text{-}\mathrm{NDFA}\Rightarrow\mathrm{DFA})$]
Sia $(T,q_0,F)$ un $\epsilon$-NDFA.

Consideriamo come insieme degli stati
\[
Q'=2^Q.
\]

La funzione di transizione del DFA è definita da
\[
T':2^Q\times A\longrightarrow 2^Q,
\]
ponendo
\[
T'(S,a):=
\bigcup_{k\ge0}
T^k(T(S,a),\epsilon).
\]

Lo stato iniziale è
\[
q_0':=
\bigcup_{k\ge0}
T^k(\{q_0\},\epsilon),
\]
mentre l'insieme degli stati finali è
\[
F':=
\{\,S\subseteq Q\mid S\cap F\neq\varnothing\,\}.
\]

È facile controllare che per ogni parola $\omega$ le computazioni ottenute tramite i rispettivi algoritmi coincidono, pertanto le decisioni dei due automi coincideranno per ogni $\omega$ e quindi decideranno lo stesso
insieme.

\end{itemize}

Ne segue che le tre nozioni di decidibilità coincidono.
\end{proof}
\end{prop}